\section{Limitations}

We study small models only, 0.6B and 1.7B, by design (Section 4). A larger model may judge itself better, so self-judging may decay less in it; the main finding, however, concerns the judge rather than the learner, and the 1.7B did not escape the collapse on $\tau^2$-bench (5.5).

All users are simulated: our own on WebShop and Gemini as the customer on $\tau^2$-bench. Real users may give more or fewer clear verdicts, and be more or less polite. We also designed the WebShop user, and the live test of the sense (5.4) uses its reactions; against this, no version of the sense trained on shopping or on our users' words, and Experiment 5 uses a customer we did not design.

Each question rests on one benchmark, in a reduced setting: WebShop's small setting (1,000 products and generated requests) for Experiments 1-4, and $\tau^2$-bench for Experiment 5. Our WebShop scores are not comparable to published ones, and the findings may depend on the task: nearly all that a WebShop life learns is to finish the purchase (5), so the results may not carry to tasks where the skill is subtler. On $\tau^2$-bench we ran no life without learning, so the falls in Experiment 5 are measured against each life's own day 0. Seeds and exams are small: 3-5 seeds for most conditions (one for the year-long \emph{oracle} and the pilot), 50 WebShop exam tasks and 20-40 per $\tau^2$-bench domain, so several gaps are significant only by a rank test and moves under about 10 points on $\tau^2$-bench are within noise. The central result pools 97 lives (5.6), and Appendix B reports every test, with Holm correction for the primary claims.

The sense changed between experiments: v1 in Experiment 2, v2 in Experiment 3 and v3 in Experiments 4 and 5, and v3 was built after Experiment 3 exposed v2's failure. Differences between experiments therefore mix the sense version with the user and the sleep settings, and v3 is not blind to the WebShop user's failure: its labelling guide includes a worked example almost word for word like Example 1, and the guide's rows on polite unhappiness and on thanks for the effort follow the behaviours of our simulated user. The live test (5.4) and Experiment 4 are therefore not independent of that failure; the only test on a user we did not design is Experiment 5. Within each experiment, the sense is compared with the \emph{oracle}, and on WebShop with \emph{both}, under identical conditions, and the live test compares v2 and v3 on the same reactions.

The \emph{oracle} differs from the other judges in more than the judge: it gets no retries and so no preference pairs, while every other judge does (3.1, 4). Part of the gap between the \emph{oracle} and the others may therefore come from the retries and the DPO term, and we ran neither an \emph{oracle} with retries nor a self-judge without them. Each judge is the simplest version of a published family that fits the nightly loop, not its original implementation (3.1), so the original methods may do better. $\tau^2$-bench ran with Gemini 3 Flash as customer and graders instead of the published GPT models, so its scores are comparable to published ones, not identical. Finally, no fix for forgetting or for the late fade was tried, and the cause of the fade is unknown.
